
%  ___                              _        _   _             
% | _ \___ _ __ _ _ ___ ___ ___ _ _| |_ __ _| |_(_)___ _ _  ___
% |   / -_) '_ \ '_/ -_)_-</ -_) ' \  _/ _` |  _| / _ \ ' \(_-<
% |_|_\___| .__/_| \___/__/\___|_||_\__\__,_|\__|_\___/_||_/__/
%         |_|                                                  

\chapter{Representations}

\noindent \textbf{Representations} are maps from a Lie group or a Lie Algebra 
to a vector space of matrices. The representations of a Lie group preserve the 
group properties and the representations of a Lie algebra preserve the Lie 
bracket. Formally:

\begin{definition}[Lie Group Representation]
    A representation of a Lie group $G$ is a map $\Pi$ from any Lie group element
    $g\in{G}$ to a linear transformation $\Pi(g)$ of some vector space $V$ such
    that group properties are preserved.
\begin{equation}
    g \underbrace{\longrightarrow~}_{\Pi} \Pi(g) \tag{Sch:3.65}
\end{equation}
\end{definition}

\begin{remark}
\noindent Preserving group properties means:
\begin{description}[noitemsep]
    \item[Multiplication is preserved:] $\Pi(g_1 g_2) = \Pi(g_1) \Pi(g_2)$ 
    \item[Group identity maps to matrix identity:] $\Pi(e) = I$
    \item[Inverses are preserved:] $\Pi(g^{-1}) = \Pi(g)^{-1}$
\end{description}
\end{remark}

\begin{definition}[Lie Algebra Representation]
    A representation of a Lie algebra $\mathfrak{g}$ is a map $\pi$ from any 
    Lie algebra element $X \in \mathfrak{g}$ to a linear transformation 
    $\pi(X)$ of some vector space $V$ such that the Lie bracket is preserved:
    \begin{equation}
        \pi([X, Y]) = [\pi(X), \pi(Y)] % \tag{Sch:3.66}
    \end{equation}
    where the bracket on the right is the standard matrix commutator: $[\pi(X), \pi(Y)] \definedas \pi(X)\pi(Y) -
  \pi(Y)\pi(X)$.
\end{definition}

\noindent When a Lie group and a Lie algebra are linked by the exponential map, every
representation of a group ``lifts'' to a representation of its algebra, and vice
versa. Figure~\ref{fig:bridge} shows the relationships amongst:

\begin{description}% [noitemsep]
    \item[The Lie algebra $\mathfrak{g}$:] e.g, \su{2}, spanned by by $T$ in the fundamental 
        and by $J$ in the adjoint representation
    \item[The Lie group $G$:] e.g., $G_{\text{sc}}=\SU{2}$, the unique simply connected
        group, or $G=\SO{3}$, the multiply connected group of rotation
        matrices in $\mathbb{R}^3$
    \item[The Carrier Space $V$:] e.g., for $\SU{2}$, $V = \mathbb{C}^2$, the 
        \firstmention{spinor space}; for $\SO{3}$, $V = \mathbb{R}^3$, the physical Euclidean space
    \item[The General Linear Group $GL(V)$:] the group of all invertible linear transformations
        of the carrier space
    \item[The General Linear Algebra $\mathfrak{gl}(V)$:] the space of all linear maps from $V$
        to $V$ that preserve the Lie bracket as the matrix commutator $[A,B]=AB-BA$
\end{description}

\begin{figure}[h]
    \centering
    \refstepcounter{figure}\label{fig:bridge}
    \addtocounter{figure}{-1}
    \begin{tikzpicture}[
        node distance=3.5cm,
        auto,
        >=Stealth
    ]
    % Nodes
    \node (g) {$\mathfrak{g}$};
    \node (G) [above of=g] {$G$};
    \node (glV) [right of=g] {$\mathfrak{gl}(V)$};
    \node (GLV) [above of=glV] {$GL(V)$};

    % Mappings (Morphisms)
    \draw[->, bend left=20] (g) to node {$\exp$} (G);
    \draw[->, dashed, bend left=20] (G) to node {$\log$} (g);

    \draw[->, bend left=20] (glV) to node {$\operatorname{Exp}$} (GLV);
    \draw[->, dashed, bend left=20] (GLV) to node {$\operatorname{Log}$} (glV);

    % The Functorial Relationship (Morphisms between Morphisms)
    \draw[->, blue, thick] (G) to node {$\Pi$} (GLV);
    \draw[->, red, thick] (g) to node [swap] {$\pi \definedas d\Pi$} (glV);

    % Legend / Labels
    \node[anchor=west] at (5, 1.7) {
        \begin{minipage}{5.5cm}
        \footnotesize
        \(\pi\) is the \firstmention{differential} of \(\Pi\). \\
        \vspace{2mm}
        \(\log \circ \exp = \operatorname{id}_{\mathfrak{g}}\) (Locally) \\
        \(\exp \circ \log = \operatorname{id}_G\) (Locally) \\
        \vspace{2mm}
        \(\operatorname{Exp} \circ \pi = \Pi \circ \exp\) \\
        \vspace{2mm}
        \(\pi\) is the \textit{linearization} of \(\Pi\).
        \end{minipage}
    };
\end{tikzpicture}
\caption{The relationship amongst Lie group and Lie algebra representations.}
\end{figure}

\section{The Two Worlds of su(2)}

\noindent In our motivating example of \su{2}, we found two distinct bases: a
$2 \times 2$ complex basis $\{T_a\}$ and a $3 \times 3$ real basis $\{J_a\}$.
To a physicist, these are not just different sizes of matrices; they represent
two different ways a group can act on the physical world.

\subsection{The Fundamental Pair (Direct Action)}
The $2 \times 2$ representation is the \firstmention{fundamental
representation}. It describes how the group \SU{2} acts on its ``native'' carrier
space $V = \mathbb{C}^2$ (the space of spinors).
\begin{itemize}[noitemsep]
    \item \textbf{Group Level}: $\Pi_{\text{fund}}(Q)$ is the $2 \times 2$ matrix itself.
    \item \textbf{Algebra Level}: $\pi_{\text{fund}}(X)$ is the $2 \times 2$ generator.
    \item \textbf{Action}: The \firstmention{direct action} $Q \ket{\xi}$, where $Q \in \SU{2}$
        and $\ket{\xi}$ is a spinor, produces a new spinor.
\end{itemize}

\subsection{The Adjoint Pair (Adjoint Action)}
The $3 \times 3$ representation is the \firstmention{adjoint representation}.
It describes how the group \SU{2} acts on its own Lie algebra
\su{2}.
\begin{itemize}[noitemsep]
    \item \textbf{Group Level}: $\Pi_{\text{adj}}(Q)$ is the $3 \times 3$ rotation matrix in \SO{3}.
    \item \textbf{Algebra Level}: $\pi_{\text{adj}}(X)$ is the $3 \times 3$ matrix of structure 
        constants $(J_a)_{bc} = \epsilon_{abc}$.
    \item \textbf{Action}: The \firstmention{adjoint action} $Q X Q^\dagger$, where $Q \in \SU{2}$ 
        and $X$ is a vector in the algebra, produces a new vector in the algebra.
\end{itemize}

\begin{remark}
Note the ``squaring'' relationship: The adjoint representation is (roughly) the
product of two fundamental representations ($Q \dots Q^\dagger$). This is why
\SO{3} is the adjoint group of \SU{2}. It is also why spinors (\SU{2}) are the
"square root" of vectors (\SO{3}), and why it takes a $4\pi$ rotation to
return a spinor to its original state, while $2\pi$ suffices for a vector.
Another view is that the relationship between the fundamental and adjoint
representations explains the \firstmention{double covering} of \SO{3} by \SU{2}.
\end{remark}

\begin{figure}[h]
    \centering
    \begin{tikzpicture}[node distance=2cm, auto, >=Stealth]
        % Fundamental World
        \node (spinor) {$\ket{\xi} \in \mathbb{C}^2$};
        \node (spinor_rot) [right=of spinor] {$Q\ket{\xi}$};
        \draw[->] (spinor) to node {Direct} (spinor_rot);
        \node at (1.5, -0.8) {\footnotesize \textbf{Fundamental (Spin-1/2)}};

        % Adjoint World
        \node (vec) [below=1.5cm of spinor] {$X \in \su{2}$};
        \node (vec_rot) [right=of vec] {$QXQ^\dagger$};
        \draw[->] (vec) to node {Adjoint} (vec_rot);
        \node at (1.5, -3.2) {\footnotesize \textbf{Adjoint (Spin-1)}};

        % Label
        \node[draw, dashed, inner sep=10pt] at (6.5, -1.5) {
            \begin{minipage}{3.5cm}
            \footnotesize
            \textbf{Topology Bridge:} \\
            \(Q \to -Q\) leaves the \\
            Adjoint action invariant, \\
            but flips the Direct action.
            \end{minipage}
        };
    \end{tikzpicture}
    \caption{Direct vs. Adjoint actions of \SU{2}.}
\end{figure}

\begin{fact}
The Adjoint representation is always available for any Lie group. It is the
"regulation" way to view a Lie group as a matrix group of rotations on its own
algebra.
\end{fact}
